\chapter{Automatische Fehlersuche mit LuaCheck}
\label{ch:quality/luacheck}

In diesem Kapitel lernst du, wie du deinen Mod mit dem Werkzeug
LuaCheck\index{luacheck@Luacheck} automatisch nach Fehlern durchsuchst.
LuaCheck lässt sich mit deinem Editor verbinden, der dich dann auf
Fehler hinweist.

\section{LuaCheck installieren}\label{luacheck-installieren}

\subsection{Windows}\label{windows}

Lade einfach \texttt{luacheck.\allowbreak{}exe} von der Release-Seite
auf GitHub\footnote{\url{https://github.com/lunarmodules/luacheck/releases}}
herunter.

\subsection{Linux}\label{linux}

Installiere zuerst LuaRocks:

\begin{Shaded}
\begin{Highlighting}[]
\NormalTok{sudo apt install luarocks}
\end{Highlighting}
\end{Shaded}

Dann installierst du LuaCheck global:

\begin{Shaded}
\begin{Highlighting}[]
\NormalTok{sudo luarocks install luacheck}
\end{Highlighting}
\end{Shaded}

Mit diesem Befehl prüfst du, ob die Installation geklappt hat:

\begin{Shaded}
\begin{Highlighting}[]
\NormalTok{luacheck {-}v}
\end{Highlighting}
\end{Shaded}

\section{LuaCheck ausführen}\label{luacheck-ausfuxfchren}

Beim ersten Durchlauf meldet LuaCheck wahrscheinlich viele Fehler, die
keine sind. Das liegt daran, dass es noch nicht eingerichtet ist.

Unter Windows öffnest du PowerShell oder Bash im Hauptordner deines
Projekts und führst
\texttt{path\textbackslash{}to\textbackslash{}luacheck.\allowbreak{}exe .\allowbreak{}}
aus.

Unter Linux führst du im Hauptordner deines Projekts
\texttt{luacheck .\allowbreak{}} aus.

\section{LuaCheck einrichten}\label{luacheck-einrichten}

Lege im Hauptordner deines Projekts eine Datei
\texttt{.\allowbreak{}luacheckrc} an. Das kann der Hauptordner deines
Spiels, Mod-Pakets oder Mods sein.

Schreib Folgendes hinein:

\begin{Shaded}
\begin{Highlighting}[]
\VariableTok{unused\_args} \OperatorTok{=} \KeywordTok{false}
\VariableTok{allow\_defined\_top} \OperatorTok{=} \KeywordTok{true}

\VariableTok{globals} \OperatorTok{=} \OperatorTok{\{}
    \StringTok{"core"}\OperatorTok{,}
\OperatorTok{\}}

\VariableTok{read\_globals} \OperatorTok{=} \OperatorTok{\{}
    \VariableTok{string} \OperatorTok{=} \OperatorTok{\{} \VariableTok{fields} \OperatorTok{=} \OperatorTok{\{} \StringTok{"split"} \OperatorTok{\}} \OperatorTok{\},}
    \VariableTok{table} \OperatorTok{=} \OperatorTok{\{} \VariableTok{fields} \OperatorTok{=} \OperatorTok{\{} \StringTok{"copy"}\OperatorTok{,} \StringTok{"getn"} \OperatorTok{\}} \OperatorTok{\},}

    \CommentTok{{-}{-} Eingebaut}
    \StringTok{"vector"}\OperatorTok{,}
    \StringTok{"ItemStack"}\OperatorTok{,}
    \StringTok{"dump"}\OperatorTok{,}
    \StringTok{"DIR\_DELIM"}\OperatorTok{,}
    \StringTok{"VoxelArea"}\OperatorTok{,}
    \StringTok{"Settings"}\OperatorTok{,}

    \CommentTok{{-}{-} Minetest Game}
    \StringTok{"default"}\OperatorTok{,}
    \StringTok{"sfinv"}\OperatorTok{,}
    \StringTok{"creative"}\OperatorTok{,}
\OperatorTok{\}}
\end{Highlighting}
\end{Shaded}

Verwendest du in älterem Code noch den
Namensraum\index{namensraum@Namensraum} \texttt{minetest}, nimm ihn
zusätzlich in \texttt{globals} auf.

Prüfe dann mit einem neuen Durchlauf, ob es funktioniert. Diesmal
solltest du deutlich weniger Fehler bekommen. Arbeite dich ab dem ersten
Fehler durch: Behebe das Problem im Code oder passe die Konfiguration
an, wenn der Code korrekt ist. Hilfe dazu gibt die folgende Liste.

\subsection{Fehlerbehebung}\label{fehlerbehebung}

\begin{itemize}
\tightlist
\item
  \textbf{accessing undefined variable foobar} -- Soll \texttt{foobar}
  eine globale Variable sein, nimm sie in
  \texttt{read\_\allowbreak{}globals} auf. Andernfalls ergänze die
  fehlenden \texttt{local}.
\item
  \textbf{setting non-standard global variable foobar} -- Soll
  \texttt{foobar} eine globale Variable sein, nimm sie in
  \texttt{globals} auf und entferne sie gegebenenfalls aus
  \texttt{read\_\allowbreak{}globals}. Andernfalls ergänze die fehlenden
  \texttt{local}.
\item
  \textbf{mutating read-only global variable `foobar'} -- Verschiebe
  \texttt{foobar} von \texttt{read\_\allowbreak{}globals} nach
  \texttt{globals} oder schreib nicht mehr in \texttt{foobar}.
\end{itemize}

\section{Im Editor verwenden}\label{im-editor-verwenden}

Es ist sehr zu empfehlen, für deinen Editor eine Erweiterung zu
installieren, die Fehler anzeigt, ohne dass du einen Befehl ausführen
musst. Für die meisten Editoren gibt es eine.

\begin{itemize}
\tightlist
\item
  \textbf{VSCode} -- Strg+P drücken und dann einfügen:
  \texttt{ext install dwenegar.\allowbreak{}vscode-luacheck}
\item
  \textbf{Sublime Text} -- Über Package Control installieren:
  SublimeLinter\footnote{\url{https://github.com/SublimeLinter/SublimeLinter}}
  und SublimeLinter-luacheck\footnote{\url{https://github.com/SublimeLinter/SublimeLinter-luacheck}}.
\end{itemize}
